%%%%%%%%%%%%%%%%%%%%%%%%%%%%%%%%%%%%%%%%%%%%%%%%%%%%%%%%%%%%%%%%%%%%%%%%%%%%%%%
%% Capítulo 3 --- Desenvolvimento e métodos
%%
%% Reúne o que os cinco relatórios de origem descreviam em aberturas de seção
%% separadas, cada um com escopo próprio. Os cinco "Escopo" colapsam aqui na
%% delimitação única da Seção 3.1.
%%
%% Ordem das subseções espelha a do Capítulo 4, para que método e resultado
%% correspondente possam ser localizados em paralelo.
%%%%%%%%%%%%%%%%%%%%%%%%%%%%%%%%%%%%%%%%%%%%%%%%%%%%%%%%%%%%%%%%%%%%%%%%%%%%%%%

\makeatletter
\@ifundefined{toprule}{%
  \newcommand{\toprule}{\hline}%
  \newcommand{\midrule}{\hline}%
  \newcommand{\bottomrule}{\hline}%
}{}
\@ifundefined{emandamento}{%
  \@ifundefined{textcolor}%
    {\newcommand{\emandamento}[1]{#1}}%
    {\newcommand{\emandamento}[1]{\textcolor{red}{#1}}}%
}{}
\makeatother

%% Ver nota sobre caminhos de figura no capitulo04.tex.
\providecommand{\kxfig}[1]{../../experiments/figures/output/#1}
\providecommand{\kxhw}[1]{../../#1}

\chapter{Desenvolvimento e métodos}
\label{cap:desenvolvimento}

\section{Delimitação}
\label{sec:dev-delimitacao}

O objeto de estudo é o Kaelix, um dispositivo de monitoramento de vibração de baixo custo
composto por acelerômetro MPU6050, microcontrolador ESP32-S3, enlace LoRa e invólucro em
ASA impresso por deposição, com detecção de anomalia embarcada e rotulagem de severidade
pela norma. As especificações de hardware, as cotas do desenho mecânico e a arquitetura de
processamento são tomadas como dadas; o que se determina é o que elas permitem medir.

A estratégia é comum a todas as questões investigadas: construir uma condição de referência
cujo resultado correto decorre de solução analítica ou de definição normativa, submetê-la à
cadeia proposta e medir o desvio introduzido pela configuração do projeto. Quando a
resposta certa é conhecida de antemão, um erro de implementação não consegue se esconder
atrás de um resultado plausível.

\emandamento{Nenhum ensaio físico foi conduzido e nenhuma medição de campo alimenta os
modelos. A construção do protótipo e seu ensaio em ambiente real são a etapa seguinte
declarada deste trabalho.} As três procedências de resultado que decorrem dessa delimitação
--- sinal sintetizado, dado real decimado e modelagem sobre especificação --- estão
definidas na Seção~\ref{sec:proc}, e cada valor do Capítulo~\ref{cap:resultados} é marcado
individualmente por elas.

\section{Modelo de sinal e cadeias de aquisição}
\label{sec:dev-sinal}

\subsection{Sinal sintetizado}

Os sinais foram sintetizados com dois mecanismos de escalas de frequência distintas.
Desbalanceamento e desalinhamento aparecem como componentes harmônicas da rotação, todas
abaixo de \SI{100}{\hertz}. Defeito de pista externa aparece como um trem de impulsos
amortecidos, repetido na taxa característica da geometria e modulando uma ressonância
estrutural. A aceleração simulada é

\begin{equation}
a(t) \;=\; \underbrace{\sum_{k=1}^{3} A_k \sin\!\left(2\pi k f_r t + \varphi_k\right)}_{\text{banda baixa}}
\;+\; \underbrace{\sum_{i} D\, e^{-\zeta_d (t-t_i)} \sin\!\left(2\pi f_{\mathrm{res}} (t-t_i)\right) u(t-t_i)}_{\text{impactos do defeito}}
\;+\; \varepsilon(t),
\label{eq:dev-sinal}
\end{equation}

\noindent
com $f_r$ a rotação, $A_k$ as amplitudes harmônicas, $D$ a amplitude do defeito, $\zeta_d$
a taxa de amortecimento do impacto, $f_{\mathrm{res}}$ a ressonância excitada,
$u(\cdot)$ o degrau unitário e $\varepsilon(t)$ ruído gaussiano branco. Os instantes de
impacto são $t_i = i/\mathrm{BPFO} + \delta_i$, em que $\delta_i$ modela o escorregamento
das esferas. As frequências características decorrem da geometria do rolamento pelas
expressões usuais de BPFO e BPFI.

\subsection{Três cadeias comparadas sobre o mesmo sinal}

A comparação isola o efeito da aquisição mantendo o sinal fixo:

\begin{enumerate}
  \item \textbf{Referência} --- amostragem a \SI{20}{\kilo\hertz}, representando um
        acelerômetro de banda larga.
  \item \textbf{Dispositivo corretamente configurado} --- decimação para
        \SI{1}{\kilo\hertz} com filtro anti-\emph{aliasing}.
  \item \textbf{Dispositivo sem filtro} --- subamostragem para \SI{1}{\kilo\hertz} sem
        filtro, representando a leitura direta do registrador quando o filtro passa-baixas
        digital interno não é configurado.
\end{enumerate}

No terceiro caso, uma componente em $f$ acima de Nyquist é observada na frequência aparente

\begin{equation}
f_{\mathrm{ap}} = \left| f - f_s \left\lfloor \frac{f}{f_s} + \frac{1}{2} \right\rceil \right|,
\label{eq:dev-alias}
\end{equation}

\noindent
com $f_s$ a taxa de amostragem e $\lfloor\cdot\rceil$ o arredondamento ao inteiro mais
próximo. É essa cadeia que corresponde ao estado atual do firmware
(Seção~\ref{sec:dev-firmware}).

\begin{figure}[!htb]
\centering
\includegraphics[width=0.88\textwidth]{\kxfig{fig8_circuito.pdf}}
\caption[Circuito de condicionamento e aquisição]{Circuito de condicionamento e aquisição. Os valores da figura são gerados
diretamente a partir das constantes do firmware, com o arquivo e a linha de origem
registrados no cabeçalho do gerador.}
\label{fig:circuito}
\end{figure}

Na comparação entre cadeias, os descritores extraídos por janela são o valor eficaz, a
curtose, o fator de crista, o pico e o pico a pico. O detector embarcado usa um conjunto
menor, idêntico no treino e no firmware: valor eficaz, curtose, fator de crista e frequência
dominante. A escolha tem consequência analisada no Capítulo~\ref{cap:resultados}: os
descritores descrevem energia, impulsividade e a linha espectral mais intensa, e nenhum
captura a relação entre harmônicas ou a fase.

\subsection{Dado real decimado}

Para as taxas de detecção por modo de falha, o sinal não é sintetizado. Utilizaram-se as
gravações do conjunto público MAFAULDA, obtidas em bancada instrumentada a
\SI{50}{\kilo\hertz} e decimadas em \emph{software} para \SI{1}{\kilo\hertz}, a banda do
dispositivo. %% TODO: citar o MAFAULDA; a chave ainda não existe em referencias.bib.
O conjunto soma \num{1951} arquivos em seis classes --- operação normal, desbalanceamento,
desalinhamento horizontal e vertical e defeito de rolamento nas posições \emph{overhang} e
\emph{underhang} ---, segmentados em \num{17559} janelas de \num{512} amostras. Apenas
\num{49} arquivos são de operação normal, e é deles que sai todo o treino, já que o detector
aprende somente a operação normal. A bancada está abaixo do escopo de potência e altura de
eixo da norma, o que a torna adequada como fonte de sinal para estudar a cadeia de medição e
inadequada como caso de aplicação normativa.

O limiar do detector é calibrado pelo quantil do escore sobre operação normal, com falso
alarme alvo de \SI{1}{\percent}. \emandamento{Os resultados desse treino não têm artefato
versionado próprio --- arquivo de resultados ou registro de execução ---: estão registrados
no histórico de alterações do repositório, e reproduzi-los exige executar o treino de novo.}

A validação emprega partição por ensaio, e não por janela, em cinco partições com os
arquivos de origem como grupos. A justificativa é o próprio
resultado de vazamento apresentado na Seção~\ref{sec:res-estimacao}: janelas do mesmo
ensaio compartilham a assinatura da montagem, e dividi-las aleatoriamente mede a
capacidade de identificar o ensaio, não a falha. Nos ensaios de rolamento, os resultados
são adicionalmente estratificados pelo nível de massa de desbalanceamento sobreposta,
porque o agregado confunde os dois efeitos.

\section{Conversão normativa e estágio de decisão}
\label{sec:dev-estimacao}

A norma define as zonas de severidade em velocidade eficaz de banda larga, e o acelerômetro
entrega aceleração \cite{iso208163_2022}. A conversão foi testada contra uma senoide cuja
integral é conhecida analiticamente, comparando cinco variantes: integração no tempo com e
sem filtro passa-alta, integração no domínio da frequência restrita à máscara normativa, e
esta última e a integração no tempo sem filtro repetidas sob viés contínuo contido na
especificação do sensor. O método correto
está estabelecido na literatura \cite{brandt2014}; o que se mede aqui é o erro das
alternativas sobre um valor de referência exato.

O estágio de decisão emprega detecção não supervisionada por Isolation Forest
\cite{liu2008}, treinada exclusivamente com dados de operação normal --- o cenário de uma
máquina sem histórico rotulado. Três aspectos foram submetidos a variação controlada: a
política de limiar, pela varredura do parâmetro de contaminação contra a fração de operação
normal marcada; a métrica, pela comparação entre área sob a curva ROC global e área parcial
restrita à região de falso alarme baixo \cite{koizumi2020}; e o protocolo de validação,
pela comparação entre partição aleatória por janela e partição por ensaio.

\section{Projeto mecânico e térmico}
\label{sec:dev-mecanica}

O invólucro foi analisado a partir das cotas do desenho técnico, por fórmulas fechadas de
resistência dos materiais. A rigidez do caminho de transmissão entre a superfície do motor
e o sensor foi modelada como associação em série de três elos --- a base como placa
circular engastada com carga central, o ressalto de centragem em compressão axial e o corpo
como viga tubular engastada:

\begin{equation}
k_{\mathrm{base}} = \frac{16\pi \mathcal{D}}{a^2},
\qquad
k_{\mathrm{res}} = \frac{E A}{L},
\qquad
k_{\mathrm{corpo}} = \frac{3 E I}{L^3},
\qquad
\frac{1}{k} = \sum_i \frac{1}{k_i},
\label{eq:dev-rigidez}
\end{equation}

\noindent
com $\mathcal{D} = E h^3 / [12(1-\nu^2)]$ a rigidez flexional da placa e
$I = [\ell^4 - (\ell-2t)^4]/12$ o momento de inércia da seção tubular quadrada. A
frequência de montagem e a transmissibilidade de base decorrem do modelo de um grau de
liberdade.

A análise térmica considera o invólucro fechado em regime permanente, com o calor entrando
pela carcaça do motor e o autoaquecimento do circuito como fonte interna. O caminho
condutivo é modelado como resistência térmica em série, e o limite de operação é
determinado pelo primeiro componente a atingir sua temperatura máxima admissível.

\emandamento{As condições de contorno são idealizadas: a rigidez real depende de pré-carga
de parafusos, planicidade das faces impressas e orientação das camadas, e a análise térmica
não modela resistência de contato, convecção forçada pelo ventilador do motor nem radiação
de superfícies próximas.}

\section{Geração e verificação do hardware}
\label{sec:dev-hardware}

Nenhum arquivo derivado é editado à mão. O esquemático, a placa e o sólido são gerados por
\emph{script} a partir de uma fonte única, e cada etapa tem verificação automática
associada, conforme a Tabela~\ref{tab:dev-cadeia}.

\begin{table}[!htb]
\centering
\caption{Cadeia de geração do hardware e verificação automática de cada etapa.}
\label{tab:dev-cadeia}
\small
\begin{tabular}{@{}lll@{}}
\toprule
Etapa & Saída & Verificação automática \\
\midrule
Esquemático & \texttt{.kicad\_sch}, netclasses & ERC com aviso tratado como erro \\
Placa       & \texttt{.kicad\_pcb}             & DRC com zonas preenchidas; colisão \\
Sólido      & \texttt{.step} da placa e do invólucro & modelo por \emph{footprint}; interferência \\
Modal       & frequências e formas             & convergência de malha em três refinos \\
Térmica     & campo e transiente               & balanço fechado contra modelo concentrado \\
\bottomrule
\end{tabular}
\end{table}

\begin{figure}[!htb]
\centering
\begin{minipage}{0.24\textwidth}\centering
\includegraphics[width=\linewidth]{\kxhw{hardware/pcb/v1/render_top.png}}\\[1pt]
{\footnotesize v1, superior}
\end{minipage}\hfill
\begin{minipage}{0.24\textwidth}\centering
\includegraphics[width=\linewidth]{\kxhw{hardware/pcb/v1/render_bottom.png}}\\[1pt]
{\footnotesize v1, inferior}
\end{minipage}\hfill
\begin{minipage}{0.24\textwidth}\centering
\includegraphics[width=\linewidth]{\kxhw{hardware/pcb/v2/render_top.png}}\\[1pt]
{\footnotesize v2, superior}
\end{minipage}\hfill
\begin{minipage}{0.24\textwidth}\centering
\includegraphics[width=\linewidth]{\kxhw{hardware/pcb/v2/render_bottom.png}}\\[1pt]
{\footnotesize v2, inferior}
\end{minipage}
\caption[Placas de circuito impresso geradas, versões v1 e v2]{Placas geradas nas duas versões. Contorno retangular com quatro recortes de canto
para os apoios do invólucro; duas camadas, com plano de terra em ambas.}
\label{fig:placas}
\end{figure}

A regeneração a partir de diretório limpo reproduz os arquivos de esquemático, placa e
símbolos byte a byte nas duas versões; o roteador é determinístico. Essa propriedade é o
que torna a verificação significativa: um resultado que não se reproduz não distingue
correção de coincidência.

As regras de fabricação adotadas são traço de \SI{0,20}{\milli\meter}, folga de
\SI{0,15}{\milli\meter} e via de \SI{0,60}{\milli\meter} com furo de
\SI{0,30}{\milli\meter}, contra mínimos declarados de \SI{0,15}{\milli\meter} para traço e
folga. \emandamento{Os defeitos de geração encontrados e corrigidos, bem como os que
permanecem abertos, constam do Apêndice~\ref{ape:hardware}.}

\begin{figure}[!htb]
\centering
\includegraphics[width=0.94\textwidth]{\kxfig{fig9_pdn.pdf}}
\caption[Rede de distribuição de alimentação]{Rede de distribuição de alimentação: impedância vista do pino de cada carga,
composição da indutância de laço e comprimento roteado contra distância em linha reta.}
\label{fig:pdn}
\end{figure}

\section{Firmware e estado de implementação}
\label{sec:dev-firmware}

O firmware é escrito em C++ para o ESP32-S3, com disciplina de sistemas críticos: caminhos
não implementados devolvem um estado de erro explícito e propagam-no, em vez de retornar
sucesso ou falha silenciosa. A justificativa dessa escolha e sua correspondência com a
prática de sistemas de alta criticidade constam do Apêndice~\ref{ape:praticas}.

\emandamento{Dois caminhos permanecem não implementados e são decisivos para a validade da
medição em campo: a inicialização do acelerômetro --- que inclui a conferência do registro
de identidade, o fundo de escala, o divisor de taxa de amostragem e, decisivamente, a
configuração do filtro passa-baixas interno abaixo de $f_s/2$ --- e o segundo caminho de
aquisição. Enquanto o filtro não for configurado, o dispositivo opera na terceira cadeia da
Seção~\ref{sec:dev-sinal}, aquela em que todo conteúdo acima de \SI{500}{\hertz} dobra para
dentro da banda de análise.}

\section{Comunicação e recepção}
\label{sec:dev-comunicacao}

A cadeia vai do motor ao manutentor em dez elos. Os parâmetros do enlace LoRa foram fixados
explicitamente --- fator de espalhamento, largura de banda e taxa de codificação ---, em vez
de herdados do padrão da biblioteca. A razão é de rastreabilidade: parâmetros herdados
significam que uma atualização de biblioteca altera o consumo do produto sem que a mudança
seja detectável.

O receptor assume três responsabilidades que o dispositivo não pode ter: carimbar o tempo,
já que o contador de partida fornece ordem e não instante; detectar lacuna na sequência,
porque perda silenciosa é indistinguível de máquina parada; e decidir quando abrir alerta.
A regra de disparo é implementada no receptor, e não no dispositivo, porque alterá-la no
dispositivo exigiria regravar cada aparelho em campo. A verificação de integridade é
simétrica: o CRC do receptor é comparado contra o do firmware compilado, em \num{204} casos
de teste.

\begin{figure}[!htb]
\centering
\includegraphics[width=\textwidth]{\kxfig{fig6_cadeia.pdf}}
\caption[Cadeia do sistema e estado de implementação de cada elo]{Cadeia do sistema e estado de implementação de cada elo. \emandamento{``Testado''
refere-se a verificação no \emph{host}: nenhum elo foi exercitado com rádio ou sensor
reais.}}
\label{fig:cadeia}
\end{figure}

\begin{figure}[!htb]
\centering
\includegraphics[width=0.92\textwidth]{\kxfig{fig5_topologia.pdf}}
\caption[Comparação entre arquiteturas topológicas]{Comparação entre arquiteturas topológicas. O enlace decide autonomia tanto quanto
decide alcance.}
\label{fig:topologia}
\end{figure}

\emandamento{A cadeia foi exercitada injetando os efeitos que o meio impõe --- colisão,
perda e corrupção de bit. O que a simulação não cobre é a propagação: alcance e margem de
enlace dependem de medição em campo, e nenhum elo foi exercitado com rádio ou sensor
reais.}
